%% toc: P15 164 164.11
\chapter{Normal Forms in the Two Model Languages of Systematic Power}
\label{sec:p15-w02c-p15-028-model-languages}

Hempel and Oppenheim introduce minimal and maximal sentences through
two small model languages before they define systematic power
\cite[pp.~166--169]{HempelOppenheim1948Explanation}. The first has one
individual constant and two unary predicates; the second adds a second
constant and a binary predicate. Their computations fix what the
general constructions of \cref{sec:p15-systematic-power} amount to in
the smallest cases, and one of them marks where the notion of minimal
content stops. This chapter states every one of those computations,
verifies each in the annex, and connects them to the content measure
and to the systematic-power example of
\cref{thm:p15-power-source-examples}.

\paragraph{Inputs and conclusions.}
The input is the finite ground vocabulary of
\cref{def:p15-power-language}, specialized to the two source languages
\(L_1\) and \(L_2\) of \cref{def:p15-w02c-p15-028-languages}, together
with the state descriptions \(D_v\), the minimal sentences \(M_v\), the
truth set \([U]\) and the content family \(\mathcal C(U)\) of that
definition. The conclusions are:
the complete lists of minimal and maximal sentences of \(L_1\); the
minimal and maximal normal forms of \(P(a)\), \(Q(a)\),
\(P(a)\to Q(a)\), a valid sentence and a contradictory sentence, with
their constituent counts \(m_U\) and \(l_U\)
(\cref{thm:p15-w02c-p15-028-l1-forms}); the eight atoms and
\(256\) minimal sentences of \(L_2\), the identification of the
source's sample minimal sentence, and the \(128\)-member minimal normal
forms of \(P(a)\) and \(Q(a)\) in \(L_2\)
(\cref{thm:p15-w02c-p15-028-l2-forms}); and the fact that the minimal
sentence \(P(a)\lor Q(a)\) has a consequence that is neither valid nor
equivalent to any singular sentence
(\cref{thm:p15-w02c-p15-028-nonsingular}).

\section{The two languages and the constituent counts}
\label{sec:p15-w02c-p15-028-counts}

\begin{hbconvention}{def:p15-w02c-p15-028-transcription}{The source's logical notation}
The source writes negation, disjunction, conjunction and the
conditional as \(\sim\), \(\mathrm v\), \(\cdot\) and \(\supset\), the
universal quantifier as \((x)\) and the existential quantifier as
\((Ex)\). This chapter writes them \(\neg\), \(\lor\), \(\land\),
\(\to\), \(\forall x\) and \(\exists x\), as
\cref{sec:p15-systematic-power} does. Predicate letters, individual
constants, the names \(L_1,L_2\) and the counts \(m,m_U,l_U\) are the
source's own.

The source calls a sentence formally true (formally false) when it
(its denial) is provable by the rules of the lower functional calculus,
and calls two sentences equivalent when each is derivable from the
other \cite[p.~158, (7.1a)]{HempelOppenheim1948Explanation}. By
\cref{thm:fol-soundness} and \cref{thm:fol-completeness},
these coincide with validity, unsatisfiability and logical equivalence
over structures with nonempty domains. The source's universe of
discourse is not restricted to the named objects
\cite[p.~158, (7.2b)]{HempelOppenheim1948Explanation}, in agreement
with \cref{def:p15-power-language}. A sentence is singular when it
contains no variables \cite[p.~158, (7.1b)]{HempelOppenheim1948Explanation}.
\end{hbconvention}

\begin{hbdefinition}{def:p15-w02c-p15-028-languages}{The model languages \(L_1\) and \(L_2\)}
\(L_1\) is the language of \cref{def:p15-power-language} whose only
individual constant is \(a\) and whose only predicates are the unary
\(P\) and \(Q\). Its ground atoms are listed as
\[
 A_1=P(a),\qquad A_2=Q(a),
\]
so \(n=2\). A ground assignment \(v=(v_1,v_2)\) is written \(v_1v_2\).

\(L_2\) adds to the vocabulary of \(L_1\) a second constant \(b\) and a
binary predicate \(R\). Its ground atoms are listed as
\[
\begin{gathered}
 A_1=P(a),\ A_2=P(b),\ A_3=Q(a),\ A_4=Q(b),\\
 A_5=R(a,a),\ A_6=R(a,b),\ A_7=R(b,a),\ A_8=R(b,b),
\end{gathered}
\]
so \(n=8\), and an assignment is written \(v_1v_2\cdots v_8\).
In both languages, \(D_v\) and \(M_v\) are formed in this atom order.
\end{hbdefinition}

\begin{hbdefinition}{def:p15-w02c-p15-028-counts}{The source's constituent counts}
In a language with \(n\) ground atoms, put \(m=2^n\). For a singular
sentence \(U\), let \(m_U\) be the number of minimal sentences in the
minimal normal form of \(U\), and \(l_U\) the number of maximal
sentences in its maximal normal form
\cite[p.~169]{HempelOppenheim1948Explanation}. In the notation of
\cref{def:p15-power-language},
\[
 m_U=|\mathcal C(U)|,\qquad l_U=|[U]|,
\]
and by \cref{thm:p15-power-normal-forms}
\[
 l_U=m-m_U .
\]
The uniform content of \cref{def:p15-power-uniform-measures} is
\(g_1(U)=m_U/m\), which is the source's formula (9.1)
\cite[p.~169]{HempelOppenheim1948Explanation}.
The minimal normal form of a valid sentence is the vacuous conjunction
and the maximal normal form of a contradictory sentence is the vacuous
disjunction \cite[pp.~166, 168]{HempelOppenheim1948Explanation}; these
are the empty conjunction \(\top\) and the empty disjunction \(\bot\)
of \cref{def:p15-power-language}.
\end{hbdefinition}

The counts depend on the language: the same sentence \(P(a)\) has
\(m_U=2\) in \(L_1\) and \(m_U=128\) in \(L_2\), while its uniform
content \(g_1\) is \(1/2\) in both.

\section{The one-name language}
\label{sec:p15-w02c-p15-028-l1}

\begin{hbtheorem}{thm:p15-w02c-p15-028-l1-forms}{proved=annex:p15-w02c-p15-028-l1-forms}{The normal forms and counts in \(L_1\)}
In \(L_1\), \(m=4\). The four minimal sentences are
\[
\begin{aligned}
 M_{00}&=P(a)\lor Q(a), & M_{01}&=P(a)\lor\neg Q(a),\\
 M_{10}&=\neg P(a)\lor Q(a), & M_{11}&=\neg P(a)\lor\neg Q(a),
\end{aligned}
\]
and the four maximal sentences (state descriptions) are
\[
\begin{aligned}
 D_{11}&=P(a)\land Q(a), & D_{10}&=P(a)\land\neg Q(a),\\
 D_{01}&=\neg P(a)\land Q(a), & D_{00}&=\neg P(a)\land\neg Q(a).
\end{aligned}
\]
These are the source's lists, in the source's order
\cite[pp.~166, 168]{HempelOppenheim1948Explanation}.
For the sentences \(U\) below, the minimal and maximal normal forms
and the counts are:
\begin{center}
\small
\begin{tabular}{@{}lllcc@{}}
 \(U\) & minimal form & maximal form & \(m_U\) & \(l_U\)\\
 \hline
 \(P(a)\) & \(M_{00}\land M_{01}\) & \(D_{11}\lor D_{10}\) & \(2\) & \(2\)\\
 \(Q(a)\) & \(M_{00}\land M_{10}\) & \(D_{11}\lor D_{01}\) & \(2\) & \(2\)\\
 \(P(a)\to Q(a)\) & \(M_{10}\) & \(D_{11}\lor D_{01}\lor D_{00}\) & \(1\) & \(3\)\\
 valid \(U\) & \(\top\) (vacuous) & \(D_{11}\lor D_{10}\lor D_{01}\lor D_{00}\) & \(0\) & \(4\)\\
 contradictory \(U\) & \(M_{00}\land M_{01}\land M_{10}\land M_{11}\) & \(\bot\) (vacuous) & \(4\) & \(0\)\\
 \(P(a)\land Q(a)\) & \(M_{00}\land M_{01}\land M_{10}\) & \(D_{11}\) & \(3\) & \(1\)
\end{tabular}
\end{center}
Written out, the minimal normal form of \(P(a)\) is
\((P(a)\lor Q(a))\land(P(a)\lor\neg Q(a))\) and its maximal normal form
is \((P(a)\land Q(a))\lor(P(a)\land\neg Q(a))\); the minimal normal form
of \(P(a)\to Q(a)\) is \(\neg P(a)\lor Q(a)\) and its maximal normal
form is
\((P(a)\land Q(a))\lor(\neg P(a)\land Q(a))\lor(\neg P(a)\land\neg Q(a))\).
In every row \(l_U=m-m_U\). For the four rows \(P(a)\),
\(P(a)\to Q(a)\), valid and contradictory, the source's own examples,
\(m_U=2,1,0,4\) respectively \cite[p.~169]{HempelOppenheim1948Explanation},
so that \(g_1(U)=1/2,1/4,0,1\).

The content family of the conjunction is the union of the content
families of its conjuncts:
\(\mathcal C(P(a)\land Q(a))=\mathcal C(P(a))\cup\mathcal C(Q(a))\).
Because \(M_{00}\) is shared, the union has three members, not four;
it is the class \(K'\) of the data \(\{P(a),Q(a)\}\) in
\cref{thm:p15-power-source-examples}.
\end{hbtheorem}

For \(P(a)\) and \(P(a)\to Q(a)\), and for the valid and contradictory
cases, the table rows repeat the source's computations
\cite[p.~168]{HempelOppenheim1948Explanation}. For \(Q(a)\) the source
gives the minimal normal form \cite[p.~166]{HempelOppenheim1948Explanation};
its maximal normal form and the row for \(P(a)\land Q(a)\) are
computed here from \cref{thm:p15-power-normal-forms}.

\section{The two-name language}
\label{sec:p15-w02c-p15-028-l2}

\begin{hbtheorem}{thm:p15-w02c-p15-028-l2-forms}{proved=annex:p15-w02c-p15-028-l2-forms}{The counts and normal forms in \(L_2\)}
\(L_2\) has exactly the eight ground atoms listed in
\cref{def:p15-w02c-p15-028-languages}, hence \(m=2^8=256\) minimal
sentences and \(256\) state descriptions
\cite[p.~166]{HempelOppenheim1948Explanation}. The source's sample
minimal sentence
\[
 P(a)\lor P(b)\lor\neg Q(a)\lor Q(b)\lor R(a,a)\lor R(a,b)
       \lor\neg R(b,a)\lor\neg R(b,b)
\]
is \(M_v\) for \(v=00100011\); it is false exactly when \(Q(a)\),
\(R(b,a)\) and \(R(b,b)\) are true and the other five atoms are false.

The minimal normal forms of \(P(a)\) and of \(Q(a)\) in \(L_2\) are
\[
 P(a)\equiv\bigwedge\{M_v: v_1=0\},\qquad
 Q(a)\equiv\bigwedge\{M_v: v_3=0\},
\]
each a conjunction of \(2^7=128\) distinct minimal sentences
\cite[p.~166]{HempelOppenheim1948Explanation}. The conjuncts for
\(P(a)\) are exactly the \(L_2\) minimal sentences having \(P(a)\) as
an unnegated disjunct, and similarly for \(Q(a)\). Hence
\(m_{P(a)}=m_{Q(a)}=128\), \(l_{P(a)}=l_{Q(a)}=128\) and
\(g_1(P(a))=g_1(Q(a))=1/2\) in \(L_2\).

Each \(L_1\) minimal sentence is singular in \(L_2\) but no longer
minimal there: \(M^{L_1}_{v_1v_2}\) has the \(L_2\) minimal normal form
consisting of the \(2^6=64\) minimal sentences \(M_w\) with
\(w_1=v_1\) and \(w_3=v_2\). Replacing each of the two conjuncts of the
\(L_1\) minimal normal form of \(P(a)\) (or of \(Q(a)\)) by its \(64\)
conjuncts in \(L_2\) yields two disjoint families whose union is the
\(128\)-member \(L_2\) minimal normal form.

More generally, if a language with ground atoms \(A_1,\ldots,A_k\) is
extended to one with ground atoms \(A_1,\ldots,A_n\), \(k\le n\), then
every singular sentence \(U\) of the smaller language has the same
uniform content \(g_1(U)\) in both, while \(m_U\) and \(l_U\) are
multiplied by \(2^{n-k}\).
\end{hbtheorem}

\section{Consequences of a minimal sentence that are not singular}
\label{sec:p15-w02c-p15-028-boundary}

\Cref{thm:p15-power-normal-forms} says that a singular consequence of a
minimal sentence is either equivalent to it or valid. The source
records that the restriction to singular consequences cannot be
removed, and names a witness
\cite[p.~166]{HempelOppenheim1948Explanation}.

\begin{hbtheorem}{thm:p15-w02c-p15-028-nonsingular}{proved=annex:p15-w02c-p15-028-nonsingular}{A consequence that is not singular}
In \(L_1\), let \(E=\exists x(P(x)\lor Q(x))\). Then:
\begin{enumerate}
\item \(P(a)\lor Q(a)\entails E\), and \(E\) is derivable from
  \(P(a)\lor Q(a)\) in the first-order calculus;
\item \(E\) is not valid, so it is not formally true;
\item \(E\) is not equivalent to any singular sentence of \(L_1\); in
  the source's terms it is essentially generalized
  \cite[p.~158, (7.1c)]{HempelOppenheim1948Explanation};
\item in particular \(E\) is not equivalent to \(P(a)\lor Q(a)\).
\end{enumerate}
So the minimal sentence \(M_{00}=P(a)\lor Q(a)\) has a consequence that
is neither equivalent to it nor valid, and the clause of
\cref{thm:p15-power-normal-forms} on consequences of \(M_v\) fails once
non-singular consequences are taken into account.

By contrast, the mutual exclusiveness of contents holds without the
singular restriction: if \(v\neq w\) and a sentence \(H\) of any form
is a consequence of both \(M_v\) and \(M_w\), then \(H\) is valid
\cite[p.~166]{HempelOppenheim1948Explanation}.

The dual clause for state descriptions has the same boundary: the
consistent sentence \(D_{11}\land\forall x\,P(x)\) entails \(D_{11}\)
but is not equivalent to it.
\end{hbtheorem}

\paragraph{Limitation.}
Minimal sentences have the smallest nonzero content only among singular
sentences. The content family \(\mathcal C(U)\), the counts \(m_U\) and
\(l_U\), and the measure \(g_1\) are defined only for singular \(U\).
A quantified theory enters the systematic-power computation only through
the minimal sentences it entails, as in \cref{thm:p15-power-optimal-split}
and the ground projection of \cref{thm:p15-power-uniform-formula}.
Its non-singular consequences, such as \(E\) above, are not counted.

\paragraph{Relationships.}
The \(L_1\) table supplies the data content used by the one-name
systematic-power example of \cref{thm:p15-power-source-examples}.
How the systematic-power score of a quantified theory changes when a
second name is added is shown in \cref{thm:p15-power-language-choice}.
The equality \(g_1(P(a))=1/2\) in both \(L_1\) and \(L_2\) is an
instance of the last clause of \cref{thm:p15-w02c-p15-028-l2-forms}:
uniform content does not change when ground atoms are added, although
the constituent counts do.

\paragraph{Sources.}
Hempel and Oppenheim, Part~IV: the language \(L_1\), its minimal
sentences, the language \(L_2\) with its counts and sample minimal
sentence, the non-singular consequence of \(P(a)\lor Q(a)\), and the
minimal normal forms of \(P(a)\) and \(Q(a)\) in \(L_1\) and \(L_2\)
\cite[p.~166]{HempelOppenheim1948Explanation}; the maximal sentences
of \(L_1\) and the four pairs of normal forms
\cite[p.~168]{HempelOppenheim1948Explanation}; the counts \(m\),
\(m_U\), \(l_U\) and formula (9.1)
\cite[p.~169]{HempelOppenheim1948Explanation}; the definitions of
formal truth, singular sentence and essentially generalized sentence
and the universe of discourse
\cite[p.~158]{HempelOppenheim1948Explanation}.
